\subsection{Thom 同构}
% ===== 原书 p.9 =====
设$V$是紧空间$X$上的复向量丛。外代数（exterior algebra）$\Lambda^*(V)$\footnote{第$n$个外幂$\Lambda^n(V)$是$n$重张量积$V^{\otimes n}$的斜对称化。}于是给出$V$上一个次数为 1 的、由向量丛组成的齐次复形$\Lambda(V)$，称为外复形（exterior complex）：
\[
0\longrightarrow\pi^*\Lambda^0(V)\xrightarrow{\ \alpha\ }\pi^*\Lambda^1(V)\xrightarrow{\ \alpha\ }\cdots\xrightarrow{\ \alpha\ }\pi^*\Lambda^n(V)\longrightarrow 0,
\]
其中$\pi:V\to X$，且对一切$v\in V$与$w\in\pi^*\Lambda^k(V)_v=\Lambda^k(V)_{\pi(v)}$，有$\alpha:(v,w)\mapsto(v,v\wedge w)$。由于$\alpha^2=0$，且该复形在零截面之外正合，可见它定义一个元素$\lambda_V\in K(V)=\tilde K(V^+)$；这里$V^+=X^V=D(V)/S(V)=P(V\oplus\mathbf{1})/P(V)$是$V$的 Thom 空间\footnote{这里把 Thom 空间定义为全空间的单点紧化。若$X$只是局部紧，则须改为逐个纤维取单点紧化。}。下面陈述$K$理论的基本定理。

\begin{theorem}[Thom 同构定理]
若$V$是紧空间$X$上的复向量丛，且把$K(V)$通过提升$\pi^*:K(X)\to K(V)$看作$K(X)$-模，则由乘以$\lambda_V$给出的 Thom 映射$\varphi:K(X)\to K(V)$是同构。
\end{theorem}

该定理的证明在附录中给出梗概。若$X$只是局部紧，则复形$\Lambda(V)$没有紧支撑。然而，若$E^\bullet$是$V$上具有紧支撑的复形，则$\Lambda(V)$与$E^\bullet$的乘积确有紧支撑，Thom 同构定理仍然成立。设$s:X\to V$是零截面，注意当$X$紧时
\[
s^*\lambda_V=s^*\circ\varphi(1)=\sum_{k=0}^n(-1)^k\Lambda^k(V).
\]
\transnote{$\lambda_V$落在$K(V)$里，而$K(V)$是$V$上紧支撑$K$理论，它靠单点紧化才等同$\tilde K(V^+)$。所以零截面上的$s^*\lambda_V$是一个$K(X)$中的类，不能把$\Lambda^k(V)$直接当作$X$上的丛来读；须先沿$s$拉回。}

% ===== 原书 p.10 =====
注意$\Lambda^*(V\oplus W)\cong\Lambda^*(V)\otimes\Lambda^*(W)$，可见$\lambda_{V\oplus W}=\lambda_V\cdot\lambda_W$，从而 Thom 同构是传递的。换言之，同构$K(X)\to K(V\oplus W)$与复合$K(X)\to K(V)\to K(V\oplus W)$相同。

\begin{example}
若$X$是一点且$V=\C^n$，则得到 Thom 同构
\[
\varphi:K(\pt)\xrightarrow{\ \cong\ }K(\C^n)\cong\tilde K(S^{2n}),
\]
从而$\tilde K(S^{2n})$同构于$\Z$，并由 Bott 类（Bott class）$\lambda_{\C^n}=\lambda_n$生成。于是$K(S^{2n})\cong\Z^2$。
\end{example}

\subsection{一个周期的上同调理论}
$K$理论最有趣也最有用的特点之一，是它具有广义上同调理论的特征性质。为看出这一点，先需要几个定义。

\begin{definition}
带基点空间$X$的约化悬挂（reduced suspension）由
\[
SX=S^1\wedge X=(S^1\times X)/(S^1\vee X)
\]
给出；$X$的$n$次迭代悬挂由
\[
S^nX=\underbrace{SS\cdots S}_{n}X=S^n\wedge X
\]
给出，其中$X\vee Y=(X\times\pt)\cup(\pt\times Y)$是两个带基点空间的单点并（one point union），$X\wedge Y=X\times Y/X\vee Y$是 smash 积。
\end{definition}

\begin{definition}
对$n\ge 0$，用
\begin{align*}
K^{-n}(X,Y)&:=\tilde K(S^n(X/Y)),\\
K^{-n}(X)&:=K^{-n}(X,\varnothing)=\tilde K(S^nX^+)
\end{align*}
定义次数平移（degree shift）\footnote{用紧支撑$K$理论，等价的定义是$K^{-n}(X):=K(\R^n\times X)$。}。
\end{definition}

\begin{corollary}[Bott 周期性]
由
\[
K^{-n}(X)\xrightarrow{\ \lambda_1\BT\cdot\ }K^{-n-2}(X)
\]
给出的映射对一切$n\ge 0$都是同构，其中$\lambda_1$是生成$\tilde K(S^2)=K^{-2}(\pt)$的 Bott 类。
\end{corollary}

\begin{proof}
取$V=\C^m\times X$为平凡丛，得到 Thom 同构
\begin{align*}
K(X)&\xrightarrow{\ \cong\ }K(\C^m\times X)=\tilde K\big((\C^m\times X)^+\big)\\
&=\tilde K(S^{2m}\wedge X^+)=K^{-2m}(X),
\end{align*}
它由
\[
x\mapsto\lambda_{X\times\C^m}\cdot\pi^*x=\pi^*_{\C^m}\lambda_m\cdot\pi^*_Xx=\lambda_m\BT x
\]
给出；这里把外张量积（external tensor product）看作映射$\BT:\tilde K(X)\otimes\tilde K(Y)\to\tilde K(X\wedge Y)$。注意$\lambda_m=(\lambda_1)^m$，所要结论立得。
\end{proof}

\transnote{此处的外张量积是带帽的$\widehat{\otimes}$，不是普通张量积。它把两个约化$K$群送进$X\wedge Y$的约化$K$群，正是悬挂同构在$K$理论里的化身；Bott 类$\lambda_m$按$\BT$相乘时才有$\lambda_m=(\lambda_1)^m$。}

% ===== 原书 p.11 =====
由 Bott 周期性，可见所有偶数维群都同构于$K^0(X)=K(X)$，所有奇数维群都同构于$K^1(X)=\tilde K(SX^+)$。于是得到$\Z_2$-分次环$K^*(X)=K^0(X)\oplus K^1(X)$，它给出一个周期的上同调理论，并伴有如下正合的「环」序列：
\[
\begin{tikzcd}[column sep=3.2em,row sep=2.4em]
K^0(X,Y)\arrow[r] & K^0(X)\arrow[r] & K^0(Y)\arrow[d]\\
K^1(Y)\arrow[u] & K^1(X)\arrow[l] & K^1(X,Y)\arrow[l]
\end{tikzcd}
\]
利用这个正合列，可以得到与整上同调类似的结果。例如，下面命题的证明留给读者（或见[6, §2.5.2]）。

\begin{proposition}
若$X$是有限 CW 复形，且胞腔只出现在偶数维，则$K^0(X)\cong\Z^r$，其中$r$是胞腔数，且$K^1(X)=0$。于是$K^*(X)\cong\Z^r$。
\end{proposition}

于是，回想复射影空间$\C\PP^n$在$0,2,\ldots,2n$每一维恰有一个胞腔，可见$K^*(\C\PP^n)\cong\Z^{n+1}$。仍得到$K(S^{2n})\cong\Z^2$。下面给出$K$理论中 K\"unneth 公式的类似物（见[6, §2.7.15]）。

\begin{theorem}[K\"unneth 定理]
设$X$、$Y$是有限 CW 复形。则有自然正合列（指标取值于$\Z_2$）
\[
0\to\sum_{i+j=k}K^i(X)\otimes K^j(Y)\to K^k(X\times Y)\to\sum_{i+j=k+1}\operatorname{Tor}\big(K^i(X),K^j(Y)\big)\to 0.
\]
\end{theorem}

\begin{example}
由于圆$S^1$上所有复向量丛都平凡，有$K(S^1)\cong\Z$。于是对一点的情形得到$K^0(\pt)\cong\Z$与$K^1(\pt)=\tilde K(S^1)=0$。另一方面，对圆$S^1$，可见$K^0(S^1)\cong\Z$与$K^1(S^1)=\tilde K(S^2)\cong\Z$。于是由 K\"unneth 定理得到$K(X\times S^1)\cong K^0(X)\oplus K^1(X)=K^*(X)$\footnote{也可取$K(X\times S^1)\cong K(X\wedge S^1)\oplus K(X)\oplus K(S^1)\oplus K(\pt)\cong\tilde K(SX^+)\oplus K(X)$。}。
\end{example}

\subsection{Chern 特征}
为讨论$K$理论与上同调的关系，需要引入示性类的概念。这里使用 \v{C}ech 上同调；处理局部紧空间时取紧支撑上同调。给定底空间$X$上的复$n$-平面丛$E$，对每个$i=0,\ldots,n$为它配以 Chern 类$c_i(E)\in H^{2i}(X;\Z)$，其中$c_0(E)=1$；并把总 Chern 类（total Chern class）定义为形式和
\[
c(E)=1+c_1(E)+\cdots+c_n(E).
\]
Chern 类的构造见[16, §14]或[13, §20]。这里只需知道 Chern 类满足下面四条重要性质\footnote{为排除所有 Chern 类都为零的退化情形，还要求$c_1(\text{Hopf 丛})$是$H^2(\C\PP^n)$的典范生成元（$n\ge 1$）。这一事实连同自然性性质与乘积公式，足以完全刻画 Chern 类。}：

\noindent\textbf{性质（自然性）}。若存在复$n$-平面丛之间的丛映射$E\to F$覆盖$f:X\to Y$，其中$E$、$F$分别是$X$、$Y$上的复$n$-平面丛，则$c(E)=f^*c(F)$。

% ===== 原书 p.12 =====
\noindent\textbf{性质（乘积公式）}。若$E$、$F$分别是同一底空间$X$上的复$m$-平面丛与复$n$-平面丛，则总 Chern 类满足
\[
c(E\oplus F)=c(E)\smile c(F),\quad\text{即}\quad c_i(E\oplus F)=\sum_{k=0}^{i}c_k(E)\smile c_{i-k}(F)\quad(i\ge 0).
\]

\noindent\textbf{性质}。若把复向量丛$E$的共轭丛（conjugate bundle）记为$\overline{E}$，于是$\overline{E}$与$E$的底实空间相同、复结构相反，则$c_k(\overline{E})=(-1)^kc_k(E)$；用总 Chern 类表示即
\[
c(\overline{E})=1-c_1(E)+c_2(E)-\cdots\pm c_n(E).
\]

\noindent\textbf{性质}。若$\mathbf{n}$是平凡复$n$-平面丛，则$c(\mathbf{n})=1$。

设$X$上的复向量丛$E$分裂为复线丛的直和$E=L_1\oplus\cdots L_n$。令$x_i=c_1(L_i)\in H^2(X;\Z)$，乘积公式给出
\[
c(E)=\prod_i c(L_i)=\prod_i(1+x_i)=1+\sum_i x_i+\sum_{i<j}x_ix_j+\cdots,
\]
可见 Chern 类$c_k(E)$是诸$x_i$的第$k$个初等对称函数（elementary symmetric function）。于是$x_i$的任一对称函数都可表为 Chern 类的多项式。虽然并非每个向量丛都这样分解\footnote{不分裂的最简单例子是$T\C\PP^2$。}，但有下面这条重要定理。

\begin{theorem}[分裂原理]
若$E$是$X$上的复向量丛，则存在空间$F(E)$与映射$\pi:F(E)\to X$，使$\pi^*:H^*(X)\to H^*(F(E))$是单射，且$\pi^*E$分裂为复线丛的直和$E=L_1\oplus\cdots\oplus L_n$。
\end{theorem}

更确切地说，取$F(E)$为$E$的旗丛（flag bundle）。证明见[6, §2.7.1]或[13, §21]。分裂原理与 Chern 类的自然性结合，说明在只涉及 Chern 类的计算中，任何向量丛都可当作线丛的直和。特别地，可以定义：

\begin{definition}
复$n$-平面丛$E$的 Chern 特征由下式给出：
\begin{align*}
\ch(E)&=\sum_i e^{x_i}=n+\sum_i x_i+\frac{1}{2!}\sum_i x_i^2+\cdots\\
&=n+c_1+\frac12(c_1^2-2c_2)+\cdots
\end{align*}
\end{definition}

\transnote{分裂原理的用法是「先拉到$F(E)$上算，再拉回」。$\pi^*$单射保证$H^*(X)$里的等式可在$H^*(F(E))$中检验，所以只用 Chern 类写出的示性类，都能在$\pi^*E$按线丛算完后唯一确定。$\ch$、$\Td$都是这样定义的。}

对线丛$L_1$、$L_2$，注意$c_1(L_1\otimes L_2)=c_1(L_1)+c_1(L_2)$，于是
\[
\ch(L_1\otimes L_2)=e^{x_1+x_2}=e^{x_1}e^{x_2}=\ch(L_1)\smile\ch(L_2).
\]
由此对任意复向量丛$E$、$F$得到
\begin{align*}
\ch(E\oplus F)&=\ch(E)+\ch(F),\\
\ch(E\otimes F)&=\ch(E)\smile\ch(F).
\end{align*}
把 Chern 特征扩充到$K(X)$，便得到自然环同态
\[
\ch:K(X)\to H^{\mathrm{even}}(X;\Q)
\]
% ===== 原书 p.13 =====
即从$K(X)$到$X$的偶数维有理上同调的自然环同态。进一步，把$\ch:\tilde K(SX)\to\tilde H^{\mathrm{even}}(SX^+;\Q)$与$\alpha^{-1}$复合，其中$\alpha:H^{\mathrm{odd}}(X;\Q)\to\tilde H^{\mathrm{even}}(SX^+;\Q)$是乘以$H^1(S^1;\Z)$的典范生成元给出的悬挂同构，就能把 Chern 特征再扩充，得到自然环同态
\[
\ch:K^1(X)\to H^{\mathrm{odd}}(X;\Q).
\]
于是可把 Chern 特征看作函子$K^*(\cdot)\to H^*(\cdot;\Q)$的自然变换；若用有理$K$理论消去挠，就得到

\begin{theorem}
若$X$是有限 CW 复形，则 Chern 特征
\[
\ch:K^*(X)\otimes\Q\to H^*(X;\Q)
\]
是有理$K$理论与有理上同调之间的自然同构。
\end{theorem}

谱序列证明见[8, §2]。偶维球面的特例在[15, §3]中结合 Bott 周期性定理的证明讨论。

\subsection{上同调的 Thom 同构}
其次，设$E$是$X$上的定向实$n$-平面丛。已知存在唯一生成元$u\in H^n(E,E_0;\Z)=H^n_c(E;\Z)$，其中$H^*_c(\cdot)$是紧支撑上同调，$E_0$是去掉$E$的零截面所得的删去空间（deleted space），使得$u$在每个纤维上的限制诱导预先给定的定向。这个生成元$u$称为$E$的上同调 Thom 类（cohomology Thom class）。再设$i:X\to E$是零截面，用$u$在$X$上的限制定义$E$的 Euler 类：
\[
e(X)=i^*u\in H^n(X;\Z).
\]
若$E$是$X$上的复$n$-平面丛，则$E$的复结构在其底实$2n$-平面丛$E_\R$上诱导一个定向。于是最高 Chern 类由$c_n(E)=e(E_\R)$定义。下面陈述上同调 Thom 同构定理。

\begin{theorem}[Thom 同构]
由$\psi:x\mapsto\pi^*x\smile u$给出的 Thom 同态
\[
\psi:H^k(X)\to H^{n+k}_c(E)
\]
是同构，其中$\pi:E\to X$是投影\footnote{与$K$理论不同，上同调的 Thom 同构定理在$E$为实丛时也适用。}。
\end{theorem}

\transnote{原文把$E$的 Euler 类记作$e(X)=i^*u$，而不是$e(E)$。原因是零截面把$X$嵌入$E$，拉回后类落在$H^n(X;\Z)$上；到第 2.3 节写$e(X)=e(TX)$时沿用同一写法。读时不要误以为这里算的是底流形的 Euler 类。}

证明见[16, §10]。用 Chern 特征可以比较$K$理论与上同调的 Thom 同构。设$E$是紧空间$X$上的复$n$-平面丛，得到下图：
\[
\begin{tikzcd}[column sep=4.2em,row sep=2.6em]
K(X)\arrow[r,"\varphi"]\arrow[d,"\ch"] & K(E)\arrow[d,"\ch"]\\
H^*(X;\Q)\arrow[r,"\psi"] & H^*_c(E;\Q)
\end{tikzcd}
\]
其中竖直映射是相应的 Chern 特征，水平映射是 Thom 同构$\varphi:x\mapsto\pi^*x\cdot\lambda_E$与$\psi:x\mapsto\pi^*x\smile u$。由于该图不交换，需要引入修正因子。确切地说，令
\[
\psi\big(\mu(E)\smile\ch(x)\big)=\ch\big(\varphi(x)\big),
\]
% ===== 原书 p.14 =====
其中
\[
\mu(E)=\psi^{-1}\circ\ch\circ\varphi(1)=\psi^{-1}\ch(\lambda_E)\in H^*(X;\Q).
\]
下面把这个上同调类$\mu(E)$明确算出。

设$i:X\to E$是零截面，回想$i^*\lambda_E=\sum_i(-1)^i\Lambda^i(E)$，并注意对一切上同调类$x\in H^*(X;\Q)$有$i^*\psi(x)=x\smile e(E_\R)$。特别地，取$x=\mu(E)$，由 Chern 类的自然性得到
\begin{align*}
\mu(E)\smile e(E_\R)&=i^*\psi\big(\mu(E)\big)=i^*\ch(\lambda_E)\\
&=\ch(i^*\lambda_E)=\ch\Big(\sum_{i=0}^n(-1)^i\Lambda^i(E)\Big).
\end{align*}
于是援引分裂原理，可设$E$分解为直和$E=L_1\oplus\cdots\oplus L_n$。注意$\Lambda^i(E)\cong\bigoplus_{j_1<\cdots<j_i}L_{j_1}\otimes\cdots\otimes L_{j_i}$，得到
\begin{align*}
\mu(E)\smile e(E_\R)&=\ch\Big(\sum_{i=0}^n(-1)^i\Lambda^i(E)\Big)\\
&=\ch\Big(\prod_{i=1}^n(1-L_i)\Big)=\prod_{i=1}^n(1-e^{x_i}),
\end{align*}
其中$x_i=c_1(L_i)\in H^2(X;\Z)$，$i=1,\ldots,n$。由$e(E_\R)=c_n(E)=\prod_i x_i$得出
\[
\mu(E)=\prod_{i=1}^n\frac{1-e^{x_i}}{x_i}.
\]
为进一步化简，再引入一个示性类。

\begin{definition}
复向量丛$E$的 Todd 类由下式给出：
\[
\Td(E)=\prod_{i=1}^n\frac{x_i}{1-e^{-x_i}}=1+\frac12c_1+\frac{1}{12}\big(c_2+c_1^2\big)+\cdots
\]
\end{definition}

注意由定义，Todd 类对任意两个向量丛$E$、$F$满足$\Td(E\oplus F)=\Td(E)\smile\Td(F)$，且对平凡复$n$-平面丛$\mathbf{n}$有$\Td(\mathbf{n})=1$。于是可见，用 Todd 类表示，修正因子为
\[
\mu(E)=(-1)^n\Td(\overline{E})^{-1};
\]
这里用到 Todd 类可逆，因其常数项为 1。

\transnote{定义式的分母是$1-e^{-x_i}$，不是$1-e^{x_i}$；换成后者会得到$\Td(E)^{-1}$那一族。修正因子里$\overline{E}$的共轭与$(-1)^n$缺一不可：$\Td(\overline{E})^{-1}=\prod_i(e^{x_i}-1)/x_i$，乘上$(-1)^n$才等于$\prod_i(1-e^{x_i})/x_i$。}
